% SynoptiQ: Project Overview
% Compile with: xelatex project_overview.tex
\documentclass[11pt,letterpaper]{article}

% ── Core packages ───────────────────────────────────────────────────
\usepackage[utf8]{inputenc}
\usepackage{fontspec}
\usepackage{microtype}
\usepackage{booktabs}
\usepackage{graphicx}
\usepackage[table]{xcolor}
\usepackage{hyperref}
\usepackage{enumitem}
\usepackage{csquotes}
\usepackage[letterpaper, top=0.9in, bottom=1.0in, left=1.0in, right=1.0in]{geometry}
\usepackage{titlesec}
\usepackage{fancyhdr}
\usepackage{tcolorbox}
\tcbuselibrary{skins,breakable}
\usepackage{array}
\usepackage{tabularx}
\usepackage{ragged2e}
\usepackage{float}
\usepackage[font={small,color=stone},labelfont={bf,color=marine},labelsep=quad]{caption}

\setlist{nosep, leftmargin=*}
\frenchspacing
\setlength{\parindent}{0pt}
\setlength{\parskip}{0.55em}

% ── Typography ─────────────────────────────────────────────────────
\setmainfont{TeX Gyre Pagella}[Ligatures=TeX]
\newfontfamily\headfont{Poppins}
\newfontfamily\greekfont{FreeSerif}
\newcommand{\gk}[1]{{\greekfont\itshape #1}}
\setmonofont{Latin Modern Mono}[Scale=0.92]

% ── Palette ────────────────────────────────────────────────────────
\definecolor{marine}{HTML}{123C48}
\definecolor{marinelight}{HTML}{2E6272}
\definecolor{terracotta}{HTML}{B75B36}
\definecolor{papyrus}{HTML}{F7F2E7}
\definecolor{stone}{HTML}{6E6A5F}
\definecolor{rulegrey}{HTML}{D9D2C0}
\definecolor{inkblack}{HTML}{201D1A}
\color{inkblack}
\arrayrulecolor{marine!65!black}

\newcolumntype{P}[1]{>{\centering\arraybackslash}p{#1}}
\newcolumntype{L}[1]{>{\RaggedRight\arraybackslash}p{#1}}
\renewcommand\labelitemi{\textcolor{terracotta}{\textbullet}}

\hypersetup{
  colorlinks=true,
  linkcolor=marine,
  citecolor=marine,
  urlcolor=terracotta,
  pdftitle={SynoptiQ: A Computational Framework for the Synoptic Problem},
  pdfauthor={Abderahmane Ainouche}
}

% ── Section styling ─────────────────────────────────────────────────
\titleformat{\section}
  {\headfont\Large\bfseries\color{marine}}
  {\color{marine}\thesection}
  {0.65em}{}
  [{\vspace{2pt}\color{rulegrey}\titlerule[0.9pt]}]
\titlespacing*{\section}{0pt}{1.7em}{0.9em}

\titleformat{\subsection}
  {\headfont\large\bfseries\color{marinelight}}
  {\thesubsection}{0.6em}{}
\titlespacing*{\subsection}{0pt}{1.3em}{0.5em}

% ── Header / footer ─────────────────────────────────────────────────
\pagestyle{fancy}
\fancyhf{}
\fancyhead[L]{\footnotesize\headfont\color{stone} SynoptiQ — Project Overview}
\fancyhead[R]{\footnotesize\headfont\color{stone} Abderahmane Ainouche}
\fancyfoot[C]{\footnotesize\headfont\color{stone}\thepage}
\renewcommand{\headrulewidth}{0.4pt}
\renewcommand{\headrule}{{\color{rulegrey}\hrule height\headrulewidth}}
\fancypagestyle{firstpage}{%
  \fancyhf{}
  \fancyfoot[C]{\footnotesize\headfont\color{stone}\thepage}
  \renewcommand{\headrulewidth}{0pt}
}

% ── Helpers ─────────────────────────────────────────────────────────
\newcommand{\statcard}[2]{%
  \cellcolor{papyrus}{\begin{minipage}[c][2.05cm][c]{\linewidth}\centering
    {\headfont\bfseries\fontsize{17}{19}\selectfont\color{marine}#1}\\[3pt]
    {\footnotesize\headfont\color{stone}#2}
  \end{minipage}}%
}
\newcommand{\tblhead}{\rowcolor{marine!12}}

\begin{document}

% ── Title block ──────────────────────────────────────────────────────
\thispagestyle{firstpage}
\begin{flushleft}
{\headfont\footnotesize\bfseries\color{terracotta}\MakeUppercase{Project Overview}}\\[7pt]
{\headfont\bfseries\fontsize{22}{27}\selectfont\color{marine}%
SynoptiQ: A Computational Framework\\[4pt] for the Synoptic Problem}\\[10pt]
{\headfont\large\color{stone} Abderahmane Ainouche}
\end{flushleft}
\vspace{6pt}
{\color{rulegrey}\rule{\linewidth}{0.9pt}}
\vspace{1.3em}

% ── Summary box ─────────────────────────────────────────────────────
\begin{tcolorbox}[
  enhanced, breakable,
  colback=papyrus, colframe=papyrus,
  boxrule=0pt, arc=0pt, outer arc=0pt,
  borderline west={2.6pt}{0pt}{marine},
  left=16pt, right=16pt, top=11pt, bottom=11pt
]
{\headfont\bfseries\color{marine} In Brief}\par\vspace{5pt}
\noindent For over two centuries, scholars have debated how the Gospels of
Matthew, Mark, and Luke are literarily related---who copied from whom, and
what lost sources might lie behind the text.  \textbf{SynoptiQ} is a
multi-phase computational framework that brings modern neural methods to
bear on this question: transformer-based language models trained on Koine
Greek, neural direction-of-copying detection, editorial fatigue modeling,
Q reconstruction, and Bayesian model comparison across the major Synoptic
hypotheses.  The first phase---a domain-adapted Koine Greek language model
and a morphologically annotated corpus of the Synoptic Gospels---is
complete and released under open licenses.  Two further phases are in
active development.
\end{tcolorbox}

\vspace{1.1em}
\begin{center}
\renewcommand{\arraystretch}{1}
\setlength{\tabcolsep}{7pt}
\begin{tabular}{P{0.18\linewidth}P{0.18\linewidth}P{0.18\linewidth}P{0.18\linewidth}P{0.18\linewidth}}
\statcard{96.62\%}{POS accuracy}
& \statcard{14~MB}{Model size}
& \statcard{49{,}061}{Tokens annotated}
& \statcard{170}{Aland pericopes}
& \statcard{3}{Papers planned}
\end{tabular}
\end{center}
\vspace{0.6em}

\section{The Question}

The Synoptic Gospels---Matthew, Mark, and Luke---share extensive verbal
parallels.  Over 90\% of Mark's verses have a parallel in Matthew, and
over half appear in Luke.  The dominant explanation, the Two-Source
Hypothesis, holds that Matthew and Luke independently used Mark and a
lost sayings source called Q.  But alternatives---the Farrer--Goulder
hypothesis (Luke used Matthew directly, no Q), the Augustinian hypothesis
(Matthew wrote first, Mark abbreviated), and the Griesbach hypothesis
(two-gospel)---each have scholarly defenders.

The evidence is textual, multi-layered, and resists simple statistical
treatment.  A scribe copying a source may follow it closely in one verse
and diverge in the next.  Editorial fatigue---the tendency of a copying
author to revert to their own stylistic norms over the course of a
pericope---leaves subtle but detectable traces.  And the direction of
literary dependence between any two parallel passages is often unclear
from surface features alone.

Computational methods offer a new way forward.  But applying them
requires building the infrastructure first: a machine-readable corpus,
a language model that understands the linguistic register of the texts,
and a suite of neural tools designed specifically for textual criticism.

\textbf{SynoptiQ is that infrastructure.}

\section{The Framework}

SynoptiQ proceeds in three stages, each yielding a standalone paper.
The full pipeline is: \textit{corpus and language model $\rightarrow$
direction detection and editorial fatigue $\rightarrow$ source
reconstruction and Bayesian comparison.}

\subsection{Paper A --- Foundation: Corpus \& Language Model \hfill \textit{Complete}}

The Synoptiq Corpus is a token-aligned, morphologically annotated dataset
of the Synoptic Gospels: 49,061 tokens across 170 Aland pericopes, with
235 pairwise Needleman--Wunsch token alignments and 13-class CCAT
part-of-speech annotations.  It is released as a HuggingFace dataset with
stratified train/validation/test splits.

KoineFormer is a domain-adapted T5 encoder-decoder trained on 1.5M tokens
of Koine Greek (SBLGNT + Apostolic Fathers) with a Classical Greek replay
buffer.  Using LoRA adapters---only 3.7M of the 220M parameters are
trained---KoineFormer achieves 96.62\% POS tagging accuracy on the
SynoptiQ test set, a 28\% error reduction over the zero-shot Classical
Greek baseline.  The full adapter checkpoint is 14 MB.

\medskip
\begin{tcolorbox}[
  enhanced, breakable,
  colback=marine!6, colframe=marine!30,
  boxrule=0.5pt, arc=3pt,
  left=12pt, right=12pt, top=8pt, bottom=8pt
]
{\headfont\bfseries\color{marine} Released}\\[4pt]
\begin{tabular}{@{}L{4cm}L{5.5cm}L{3cm}@{}}
\toprule
\tblhead\textbf{Artifact} & \textbf{Location} & \textbf{License} \\
\midrule
KoineFormer model & hf.co/ainouche-abderahmane/koineformer & CC-BY-SA 4.0 \\
Synoptiq Corpus & hf.co/datasets/ainouche-abderahmane/synoptiq-corpus & CC-BY-SA 4.0 \\
Paper A manuscript & Attached (\texttt{main.tex}) & CC-BY-SA 4.0 \\
Full codebase & github.com/abderahmane-ai/SynoptiQ & CC-BY-SA 4.0 \\
\bottomrule
\end{tabular}
\end{tcolorbox}

\subsection{Paper B --- Direction Detection \& Editorial Fatigue \hfill \textit{In Development}}

The core question---``who copied from whom?''---is addressed with two
complementary neural methods.

\textbf{Direction Scorer.}  Cross-attention asymmetry between parallel
passages produces directional features fed into a 3-way classifier
(A$\rightarrow$B, B$\rightarrow$A, independent).  An adversarial gradient
reversal layer strips authorship style to isolate the copying signal.
Training uses triple-tradition pericopes where the direction is known
(Mark$\rightarrow$Matthew, Mark$\rightarrow$Luke under the consensus
view).

\textbf{Editorial Fatigue Model.}  Position-weighted consistency loss
quantifies the well-documented phenomenon that a copying author tends to
revert to their own stylistic norms over the course of a pericope:
\(L_{\text{fatigue}} = \sum w(i) \cdot D_{\text{KL}}(\text{edit}_i \,\|\,
\text{source}_i)\), where \(w(i) = \exp(-\lambda \cdot i/N)\) decays
exponentially with position.  This provides converging evidence from an
independent signal.

\subsection{Paper C --- Q Reconstruction \& Bayesian Comparison \hfill \textit{Planned}}

\textbf{Q Reconstruction.}  A Fusion-in-Decoder architecture encodes
Matthew and Luke independently, concatenates their hidden states, and
uses a decoder with cross-attention to reconstruct the shared source.
Training proceeds in two stages: first on triple tradition (Matthew +
Luke $\rightarrow$ reconstruct Mark, where we have ground truth), then
transferred to double tradition for the hypothesized Q material.

\textbf{Bayesian Model Comparison.}  Direction scorer outputs pass
through MC Dropout ($T = 20$) to produce per-pericope uncertainty
estimates ($\mu_i, \sigma^2_i$).  These feed PyMC Beta hierarchical
models representing the four major hypotheses: Two-Source (2SH),
Farrer--Goulder (FGH), Augustinian, and Griesbach.  Bridge sampling
yields Bayes factors for formal model comparison with full uncertainty
quantification.

\section{What Distinguishes This Work}

Most computational approaches to the Synoptic Problem have used
surface-level features: word frequencies, n-gram overlap, stylometric
distances.  SynoptiQ operates at a deeper level:

\begin{itemize}
\item \textbf{Contextual representations.}  Transformer hidden states
  capture syntactic and semantic patterns that surface features miss.
\item \textbf{Causal direction detection.}  Neural cross-attention
  asymmetry provides a directional signal, rather than symmetric
  similarity measures.
\item \textbf{Adversarial debiasing.}  Gradient reversal separates the
  copying signal from authorial style---critical when the same author
  appears in multiple roles.
\item \textbf{Full uncertainty quantification.}  Bayesian hierarchical
  models with bridge sampling yield Bayes factors and credible intervals,
  not just point estimates.
\item \textbf{Complete transparency.}  Model weights, dataset, training
  code, evaluation scripts, and loss curves are all public.
\end{itemize}

\section{Current Status and Timeline}

\begin{table}[H]
\centering
\begin{tabular}{lL{3.5cm}L{5cm}l}
\toprule
\tblhead\textbf{Phase} & \textbf{Component} & \textbf{Key Deliverable} & \textbf{Status} \\
\midrule
1 & Corpus \& alignment & Synoptiq Corpus (HF dataset) & \checkmark\ Complete \\
2A & KoineFormer DAPT & Domain-adapted LM (HF model) & \checkmark\ Complete \\
2B & Multi-task fine-tuning & POS, parser, pericope heads & Code ready \\
3 & Direction scorer & 3-way classifier + GRL & Designed \\
4 & Editorial fatigue & Position-weighted KL model & Designed \\
5 & Q reconstruction & FiD transfer learning & Architecture ready \\
6 & Bayesian comparison & PyMC hierarchical + BF & Specified \\
7 & Interpretability & SHAP + BERTViz + Hawkins & Specified \\
\bottomrule
\end{tabular}
\caption{Phases of the SynoptiQ project.  Phase~2A is complete and documented in the accompanying paper (\texttt{main.tex}).  Phases~2B--7 are in various stages of development.}
\end{table}

\section{Why This Matters}

The Synoptic Problem is not merely an antiquarian puzzle.  How one
resolves it affects how one reads the Gospels---which passages are
independent witnesses, which are literary reworkings, and what can be
inferred about the historical Jesus.  A computational framework that
brings rigorous, reproducible evidence to bear on these questions serves
both New Testament scholarship and the broader digital humanities.

SynoptiQ is designed to be \textit{cumulative}.  Each phase produces
standalone results, but the full value emerges from the integration:
directional evidence from Paper~B feeds the Bayesian priors of Paper~C;
the encoder from Paper~A underlies everything.  The project is open from
the start---model weights, dataset, and code are already public---so that
scholars can inspect, reproduce, and build on each stage.

\vspace{1em}
\begin{flushright}
\textit{--- Abderahmane Ainouche}\\[4pt]
{\footnotesize\headfont\color{stone}
github.com/abderahmane-ai/SynoptiQ \hspace{1.5em}
ainouche-abderahmane/synoptiq-corpus (HF)}
\end{flushright}

\end{document}
